\documentclass[12pt]{article}
\usepackage{fullpage}
\usepackage[T1]{fontenc}
\usepackage[utf8]{inputenc}
\usepackage{lmodern}
\usepackage{microtype}
\usepackage{amsmath,amssymb,amsthm}
\usepackage{graphicx}
\usepackage{hyperref}

\newtheorem{theorem}{Theorem}
\newtheorem{lemma}{Lemma}
\newtheorem{proposition}{Proposition}
\theoremstyle{definition}
\newtheorem{remark}{Remark}

\DeclareMathOperator{\Real}{Re}
\DeclareMathOperator{\Imag}{Im}
\newcommand{\jj}{\mathrm{i}}

\begin{document}

\title{A matrix-theoretic proof of the Kharitonov vertex stability criterion}
\author{}
\date{}
\maketitle

\section{Setup and statement}

For real numbers $a_i^-\le a_i^+$, $i=0,1,\dots,n-1$, let
\[
\mathcal P=\Big\{\,p(s)=s^n+a_{n-1}s^{n-1}+\cdots+a_1s+a_0:\ a_i\in[a_i^-,a_i^+]\Big\},
\]
and let $\mathcal A=\{A(p):p\in\mathcal P\}$ be the corresponding family of companion
matrices,
\[
A(p)=\begin{pmatrix}
0&1&0&\cdots&0\\
0&0&1&\cdots&0\\
\vdots&&&\ddots&\vdots\\
0&0&0&\cdots&1\\
-a_0&-a_1&-a_2&\cdots&-a_{n-1}
\end{pmatrix}.
\]
A real square matrix is \emph{Hurwitz} (stable) if all its eigenvalues lie in the open
left half-plane $\{z:\Real z<0\}$. The four Kharitonov vertex matrices
$M_0=\{A_1,A_2,A_3,A_4\}\subset\mathcal A$ are the companion matrices of the four
polynomials whose coefficient sign patterns are, with the index $i$ read modulo $4$,
\begin{equation}\label{eq:patterns}
\begin{aligned}
A_1&:\ (-,-,+,+,-,-,+,+,\dots),\qquad
A_2:\ (+,+,-,-,+,+,-,-,\dots),\\
A_3&:\ (+,-,-,+,+,-,-,+,\dots),\qquad
A_4:\ (-,+,+,-,-,+,+,-,\dots),
\end{aligned}
\end{equation}
where the symbol in position $i$ indicates that $a_i$ is taken equal to $a_i^-$ (for $-$)
or to $a_i^+$ (for $+$).

\begin{theorem}\label{thm:main}
The whole family $\mathcal A$ is Hurwitz if and only if all four matrices in $M_0$ are
Hurwitz.
\end{theorem}

The proof below is purely matrix-theoretic: it works throughout with the matrices
$A(p)$, their characteristic functions viewed as the determinants
$\chi_p(s)=\det(sI-A(p))$, the location of their eigenvalues, and continuity/argument
arguments for these determinants. It does not invoke Kharitonov's polynomial-coefficient
theorem. The two structural ingredients --- a monotone-phase (Mikhailov) lemma for a
single stable matrix and a one-dimensional segment lemma --- are exactly the ingredients
that generalise to convex polytopes of matrices (the segment lemma is the matrix Edge
Theorem in dimension one), as discussed in Remark~\ref{rem:gen}.

\smallskip
We make no nondegeneracy assumption: the bounds may satisfy $a_i^-=a_i^+$ for some or all
$i$, and the parameter $n$ may equal $1$. In such cases the value-set ``rectangle'' below
may degenerate to a segment or a point and some of the four vertex matrices may coincide;
this affects none of the arguments, which never divide by the width of an interval and rely
only on continuity, the location of eigenvalues, and the convexity of the coefficient box.

\section{The characteristic determinant on the imaginary axis}

Throughout, for $p\in\mathcal P$ we write
\begin{equation}\label{eq:chidef}
\chi_p(s):=\det\bigl(sI-A(p)\bigr).
\end{equation}

\begin{lemma}\label{lem:det}
For every $p\in\mathcal P$ and every $s\in\mathbb C$,
\[
\chi_p(s)=s^n+a_{n-1}s^{n-1}+\cdots+a_1s+a_0 .
\]
In particular $\lambda$ is an eigenvalue of $A(p)$ if and only if $\chi_p(\lambda)=0$, and
$A(p)$ has an eigenvalue on the imaginary axis if and only if there is an
$\omega\in\mathbb R$ with $\chi_p(\jj\omega)=0$.
\end{lemma}

\begin{proof}
Expand $\det(sI-A(p))$ along the first column. The matrix $sI-A(p)$ has first column
$(s,0,\dots,0,a_0)^{\!\top}$, super-diagonal entries $-1$, diagonal entries $s$, and last
row $(a_0,a_1,\dots,a_{n-2},s+a_{n-1})$. Cofactor expansion along the first column gives
two contributions. The $(1,1)$ entry $s$ multiplies the determinant of the
$(n-1)\times(n-1)$ companion-type block of $sI-A(p)$ obtained by deleting row $1$ and
column $1$; by induction on $n$ this minor equals
$s^{n-1}+a_{n-1}s^{n-2}+\cdots+a_1$. The $(n,1)$ entry $a_0$ contributes, after the sign
$(-1)^{n+1}$, the determinant of an upper-triangular matrix with all diagonal entries
$-1$ (the $(n-1)\times(n-1)$ matrix obtained by deleting row $n$ and column $1$ is upper
triangular with $-1$'s on its diagonal), which equals $(-1)^{n-1}$; thus this term equals
$(-1)^{n+1}a_0(-1)^{n-1}=a_0$. Adding,
\[
\chi_p(s)=s\bigl(s^{n-1}+a_{n-1}s^{n-2}+\cdots+a_1\bigr)+a_0
=s^n+a_{n-1}s^{n-1}+\cdots+a_1s+a_0 .
\]
The base case $n=1$ is $\det(s+a_0)=s+a_0$. Since the eigenvalues of $A(p)$ are exactly
the zeros of $\det(sI-A(p))$, the remaining assertions follow.
\end{proof}

Because of Lemma~\ref{lem:det} we may, and do, study eigenvalue location through the
scalar functions $\chi_p$, while never leaving the matrix family $\mathcal A$: each
$\chi_p$ is the determinant of a member of $\mathcal A$.

\medskip

It is convenient to split $\chi_p(\jj\omega)$ into real and imaginary parts. For
$\omega\in\mathbb R$, since $(\jj\omega)^i=\jj^{\,i}\omega^i$ and $\jj^i$ is real for even
$i$ and purely imaginary for odd $i$, we have $\chi_p(\jj\omega)=E_p(\omega)+\jj\,O_p(\omega)$
with
\begin{equation}\label{eq:EO}
E_p(\omega)=\Real\chi_p(\jj\omega)=\sum_{i\ \mathrm{even}}(-1)^{i/2}a_i\,\omega^i,\qquad
O_p(\omega)=\Imag\chi_p(\jj\omega)=\sum_{i\ \mathrm{odd}}(-1)^{(i-1)/2}a_i\,\omega^i,
\end{equation}
where for definiteness $a_n:=1$ is included in the appropriate sum according to the
parity of $n$. Crucially:
\begin{itemize}
\item $E_p(\omega)$ depends only on the coefficients $a_i$ with $i$ even;
\item $O_p(\omega)$ depends only on the coefficients $a_i$ with $i$ odd.
\end{itemize}
Moreover, for $\omega\ge0$ each power $\omega^i\ge 0$, so the sign of the contribution of
$a_i$ to $E_p$ (resp. $O_p$) is the fixed sign $(-1)^{i/2}$ (resp. $(-1)^{(i-1)/2}$),
\emph{independent of $\omega$}. These two observations drive the whole proof.

\section{The value set is a rectangle whose corners are the Kharitonov matrices}

For $\omega\ge0$ define the \emph{value set}
\[
R(\omega):=\{\chi_p(\jj\omega):p\in\mathcal P\}\subset\mathbb C .
\]

\begin{lemma}\label{lem:rect}
For every $\omega\ge0$, $R(\omega)$ is the closed axis-parallel rectangle
\[
R(\omega)=\bigl[\,E^-(\omega),E^+(\omega)\,\bigr]\times\bigl[\,O^-(\omega),O^+(\omega)\,\bigr],
\]
where $E^\pm,O^\pm$ are the extreme values of $E_p,O_p$ over $p\in\mathcal P$. Its four
corners are
\[
\chi_{A_1}(\jj\omega)=E^-+\jj O^-,\quad
\chi_{A_2}(\jj\omega)=E^++\jj O^+,\quad
\chi_{A_3}(\jj\omega)=E^++\jj O^-,\quad
\chi_{A_4}(\jj\omega)=E^-+\jj O^+,
\]
(the dependence on $\omega$ being suppressed on the right). In particular the four
corners are attained by the four Kharitonov matrices and are the same matrices for every
$\omega\ge0$.
\end{lemma}

\begin{proof}
Fix $\omega\ge0$. As $p$ ranges over $\mathcal P$, the even coefficients $a_i$
($i$ even) range independently over $[a_i^-,a_i^+]$ and the odd coefficients $a_i$
($i$ odd) range independently over $[a_i^-,a_i^+]$. By \eqref{eq:EO}, $E_p(\omega)$ is an
affine function of the even coefficients only, and $O_p(\omega)$ is an affine function of
the odd coefficients only. Since the even and odd coefficient sets are disjoint, the pair
$(E_p(\omega),O_p(\omega))$ ranges over the full Cartesian product of the ranges of its
two coordinates as $p$ varies; hence $R(\omega)$ is the stated rectangle, and the range of
each coordinate is the interval between its extreme values $E^\pm(\omega)$,
$O^\pm(\omega)$, attained at the corresponding choices of the coefficient bounds.

We identify which bounds give the extremes. Because $\omega\ge0$, by \eqref{eq:EO} the
contribution of $a_i$ ($i$ even) to $E_p(\omega)$ has the fixed sign $(-1)^{i/2}$. Hence
$E_p(\omega)$ is maximised by choosing $a_i=a_i^+$ when $(-1)^{i/2}>0$, i.e. $i\equiv0
\pmod4$, and $a_i=a_i^-$ when $(-1)^{i/2}<0$, i.e. $i\equiv2\pmod4$; it is minimised by
the opposite choices. Likewise $O_p(\omega)$ is maximised by $a_i=a_i^+$ when
$i\equiv1\pmod4$ and $a_i=a_i^-$ when $i\equiv3\pmod4$, and minimised by the opposite
choices. Now read the sign choices of $A_2$ off \eqref{eq:patterns}: its pattern is
$(+,+,-,-,\dots)$, i.e. $a_i=a_i^+$ exactly when $i\equiv0,1\pmod4$ and $a_i=a_i^-$ when
$i\equiv2,3\pmod4$. Restricting to even $i$, this means $a_i=a_i^+$ for $i\equiv0\pmod4$
and $a_i=a_i^-$ for $i\equiv2\pmod4$, which is exactly the maximiser of $E_p$; restricting
to odd $i$, it means $a_i=a_i^+$ for $i\equiv1\pmod4$ and $a_i=a_i^-$ for $i\equiv3\pmod4$,
exactly the maximiser of $O_p$. Hence $\chi_{A_2}(\jj\omega)=E^++\jj O^+$. The pattern of
$A_1$ in \eqref{eq:patterns} is the entrywise opposite, $(-,-,+,+,\dots)$, giving the
minimisers of both $E_p$ and $O_p$, so $\chi_{A_1}(\jj\omega)=E^-+\jj O^-$. The pattern of
$A_3$, namely $(+,-,-,+,\dots)$, agrees with $A_2$ on even indices (so it maximises $E_p$)
and with $A_1$ on odd indices (so it minimises $O_p$), giving
$\chi_{A_3}(\jj\omega)=E^++\jj O^-$. The pattern of $A_4$, namely $(-,+,+,-,\dots)$, agrees
with $A_1$ on even indices and with $A_2$ on odd indices, giving
$\chi_{A_4}(\jj\omega)=E^-+\jj O^+$. (That $A_1,A_3$ share their odd-index choices, $A_2,A_4$
share their odd-index choices, $A_1,A_4$ share their even-index choices, and $A_2,A_3$ share
their even-index choices is immediate from \eqref{eq:patterns}: shifting the index by $2$
exchanges $A_1\leftrightarrow A_3$ and $A_2\leftrightarrow A_4$ on a fixed parity class.)
Since all these sign choices depend only on the residue of $i$ modulo $4$, not on $\omega$,
the same four matrices realise the four corners for every $\omega\ge0$.
\end{proof}

\section{A monotone-phase (Mikhailov) lemma for a single stable matrix}

\begin{lemma}\label{lem:mik}
Let $B=A(p)$ with $p\in\mathcal P$. Then $B$ is Hurwitz if and only if the curve
$\omega\mapsto\chi_p(\jj\omega)$ satisfies
\begin{enumerate}
\item[(M1)] $\chi_p(\jj\omega)\ne0$ for all $\omega\in\mathbb R$, and
\item[(M2)] the continuous phase $\phi_p(\omega):=\arg\chi_p(\jj\omega)$ is strictly
increasing on $\mathbb R$, with total increase $\phi_p(+\infty)-\phi_p(-\infty)=n\pi$.
\end{enumerate}
Equivalently, on $[0,\infty)$ the phase increases strictly from $\phi_p(0)=0$ (recall
$\chi_p(\jj 0)=a_0>0$ for a Hurwitz $B$) by a total amount $n\pi/2$.
\end{lemma}

\begin{proof}
Let $\lambda_1,\dots,\lambda_n$ be the eigenvalues of $B$ (with multiplicity); these are
the zeros of $\chi_p$ by Lemma~\ref{lem:det}, and since $p$ is monic of degree $n$,
\[
\chi_p(\jj\omega)=\prod_{k=1}^n(\jj\omega-\lambda_k).
\]

\emph{($\Rightarrow$)} Suppose $B$ is Hurwitz, so $\Real\lambda_k<0$ for all $k$. Write
$\lambda_k=-\alpha_k+\jj\beta_k$ with $\alpha_k>0$. Then for $\omega\in\mathbb R$,
\[
\jj\omega-\lambda_k=\alpha_k+\jj(\omega-\beta_k)\ne0\quad(\text{since }\alpha_k>0),
\]
which gives (M1). A continuous branch of the argument of $\jj\omega-\lambda_k$ is
$\theta_k(\omega)=\arctan\!\big((\omega-\beta_k)/\alpha_k\big)$, well defined because the
real part $\alpha_k>0$ keeps $\jj\omega-\lambda_k$ in the open right half-plane; its
derivative is
\[
\theta_k'(\omega)=\frac{\alpha_k}{\alpha_k^2+(\omega-\beta_k)^2}>0 .
\]
Hence $\phi_p=\sum_k\theta_k$ is a continuous phase of $\chi_p(\jj\omega)$, is strictly
increasing, and
\[
\phi_p(+\infty)-\phi_p(-\infty)=\sum_{k=1}^n\Bigl(\tfrac{\pi}{2}-\bigl(-\tfrac{\pi}{2}
\bigr)\Bigr)=n\pi,
\]
which is (M2). Finally, since the $\lambda_k$ occur in conjugate pairs together with the
real ones (the coefficients are real), $\chi_p(\jj\omega)$ is real for $\omega=0$ and equals
$a_0=\prod_k(-\lambda_k)$; for a Hurwitz polynomial $a_0>0$ (the product of $n$ numbers
$-\lambda_k$ each lying in the open right half-plane is a number whose argument is the sum
of $n$ arguments each in $(-\pi/2,\pi/2)$ but which is real, hence its argument is $0$, so
it is positive). Thus $\phi_p(0)=0$ and the half-line statement follows from the symmetry
$\chi_p(-\jj\omega)=\overline{\chi_p(\jj\omega)}$, which gives $\phi_p(-\omega)=-\phi_p(\omega)$.

\emph{($\Leftarrow$)} Conversely, assume (M1)–(M2) and suppose, for contradiction, that some
eigenvalue $\lambda_{k_0}$ has $\Real\lambda_{k_0}\ge0$. Group the factors. For each
$\lambda_k$ with $\Real\lambda_k<0$ the argument $\theta_k(\omega)$ defined above increases
from $-\pi/2$ to $\pi/2$, a net change $+\pi$. For each $\lambda_k$ with
$\Real\lambda_k>0$, writing $\lambda_k=\alpha_k+\jj\beta_k$ with $\alpha_k>0$ we get
$\jj\omega-\lambda_k=-\alpha_k+\jj(\omega-\beta_k)$ in the open left half-plane, and a
continuous argument changes by $-\pi$ as $\omega$ runs over $\mathbb R$. For each
$\lambda_k$ with $\Real\lambda_k=0$, $\jj\omega-\lambda_k$ vanishes at $\omega=\Imag
\lambda_k$, contradicting (M1); hence (M1) forces $\Real\lambda_k\ne0$ for all $k$. Let
$n_-$ and $n_+$ be the numbers of eigenvalues with negative and positive real part; then
$n_-+n_+=n$ and the total phase change of $\chi_p$ over $\mathbb R$ is
$(n_--n_+)\pi$. By (M2) this equals $n\pi$, so $n_-=n$, i.e. all eigenvalues have negative
real part, contradicting $\Real\lambda_{k_0}\ge0$. Hence $B$ is Hurwitz.
\end{proof}

\section{A one-dimensional segment lemma}

For two matrices $A(p),A(q)\in\mathcal A$ say that the \emph{even segment} between them is
$\{A(r):r=(1-t)p+tq,\ t\in[0,1]\}$ when $p$ and $q$ \emph{differ only in even-indexed
coefficients} (so $O_p\equiv O_q$); the \emph{odd segment} is defined analogously when
they differ only in odd-indexed coefficients (so $E_p\equiv E_q$). Note that for such a
convex combination $r$, the coefficient vector lies in the box, so $A(r)\in\mathcal A$.

\begin{lemma}\label{lem:seg}
If both endpoints $A(p),A(q)$ of an even segment are Hurwitz, then every matrix on that
even segment is Hurwitz. The same holds for odd segments.
\end{lemma}

\begin{proof}
We prove the even case; the odd case is identical with the roles of $E$ and $O$ exchanged.
Set $r_t:=(1-t)p+tq$, so $O_{r_t}\equiv O:=O_p=O_q$ for all $t$, while
$E_{r_t}=(1-t)E_p+tE_q$.

\emph{Step 1: no member of the segment has an imaginary eigenvalue.} Fix $t\in[0,1]$ and
$\omega\in\mathbb R$. We must show $\chi_{r_t}(\jj\omega)\ne0$, i.e.
$E_{r_t}(\omega)+\jj O(\omega)\ne0$. If $O(\omega)\ne0$ we are done. Suppose
$O(\omega)=0$. By symmetry $O$ is odd in $\omega$, so it suffices to treat $\omega\ge0$.
At such $\omega$,
\[
\chi_p(\jj\omega)=E_p(\omega),\qquad \chi_q(\jj\omega)=E_q(\omega)
\]
are both real and, since $A(p),A(q)$ are Hurwitz, both \emph{nonzero} (Lemma~\ref{lem:mik}
(M1)). We claim $E_p(\omega)$ and $E_q(\omega)$ have the \emph{same sign}; then
$E_{r_t}(\omega)=(1-t)E_p(\omega)+tE_q(\omega)$, a convex combination of two nonzero reals
of the same sign, is itself nonzero, giving $\chi_{r_t}(\jj\omega)\ne0$.

To prove the claim, recall $O_p\equiv O_q\equiv O$, so $p$ and $q$ have exactly the same
set of real zeros of $O$ on $[0,\infty)$. List the zeros of $O$ on $(0,\infty)$ together
with $0$ as
\[
0=\omega_0<\omega_1<\omega_2<\cdots ;
\]
($O(0)=0$ because $O$ has only odd powers of $\omega$, see \eqref{eq:EO}). By
Lemma~\ref{lem:mik}, for a Hurwitz $A(p)$ the phase $\phi_p$ is continuous, strictly
increasing on $[0,\infty)$, and $\phi_p(0)=0$; moreover $\phi_p(\omega)$ is an integer
multiple of $\pi$ exactly when $\chi_p(\jj\omega)$ is a nonzero real, i.e. exactly when
$O(\omega)=0$, i.e. at the points $\omega_j$. Because $\phi_p$ is strictly increasing and
$\chi_p(\jj\omega)\ne0$ throughout, $\phi_p$ passes through the values
$0,\pi,2\pi,\dots$ \emph{in order} and cannot skip any, so necessarily
\[
\phi_p(\omega_j)=j\pi\qquad(j=0,1,2,\dots),
\]
and likewise $\phi_q(\omega_j)=j\pi$, \emph{using only that the $\omega_j$ are the common
zeros of $O$ and that $\phi_p,\phi_q$ are strictly increasing with $\phi(0)=0$.}
[Indeed, strict monotonicity together with the fact that the multiples of $\pi$ attained
by $\phi_p$ occur precisely at the $\omega_j$ forces $\phi_p(\omega_j)$ to be the $j$-th
such multiple counted from $\phi_p(\omega_0)=0$, namely $j\pi$.] Consequently
\[
\operatorname{sign}E_p(\omega_j)=\operatorname{sign}\bigl(\cos\phi_p(\omega_j)\bigr)
=\cos(j\pi)=(-1)^j=\operatorname{sign}E_q(\omega_j),
\]
using $\chi_p(\jj\omega_j)=E_p(\omega_j)=|\chi_p(\jj\omega_j)|\cos\phi_p(\omega_j)$ with
$O(\omega_j)=0$. Thus $E_p$ and $E_q$ have the same sign at every common zero $\omega_j$
of $O$ in $[0,\infty)$, proving the claim and Step~1.

\emph{Step 2: stability is preserved along the segment.} The map $t\mapsto A(r_t)$ is
continuous, the endpoint $A(r_0)=A(p)$ is Hurwitz, and eigenvalues of a matrix depend
continuously on its entries (the eigenvalues are the zeros of $\det(sI-A(r_t))$, whose
coefficients are continuous in $t$; the zeros of a monic polynomial depend continuously on
its coefficients). Suppose some $A(r_{t^*})$ were not Hurwitz. Let
\[
t_0:=\sup\{\tau\in[0,1]:A(r_t)\text{ is Hurwitz for all }t\in[0,\tau]\}.
\]
Then $t_0\le t^*\le 1$, and $A(r_{t_0})$ is a limit of Hurwitz matrices, so all its
eigenvalues satisfy $\Real\lambda\le0$; by Step~1 none lies on the imaginary axis, hence
all satisfy $\Real\lambda<0$, i.e. $A(r_{t_0})$ is Hurwitz. If $t_0<1$, then by continuity
$A(r_t)$ stays Hurwitz for $t$ in a neighbourhood of $t_0$ (eigenvalues vary continuously
and remain in the open left half-plane), contradicting the definition of $t_0$ as a
supremum. Hence $t_0=1$ and $A(r_t)$ is Hurwitz for all $t\in[0,1]$.
\end{proof}

\section{Zero exclusion: the rectangle never contains the origin}

\begin{lemma}\label{lem:zero}
Suppose all four matrices $A_1,A_2,A_3,A_4$ are Hurwitz. Then $0\notin R(\omega)$ for all
$\omega\ge0$.
\end{lemma}

\begin{proof}
By Lemma~\ref{lem:rect}, $0\in R(\omega)$ if and only if
\begin{equation}\label{eq:inrect}
E^-(\omega)\le0\le E^+(\omega)\quad\text{and}\quad O^-(\omega)\le0\le O^+(\omega).
\end{equation}

\emph{The four edges are origin-free.} Consider the bottom edge of the rectangle, the
horizontal segment from the corner $\chi_{A_1}(\jj\omega)=E^-+\jj O^-$ to
$\chi_{A_3}(\jj\omega)=E^++\jj O^-$. By Lemma~\ref{lem:rect}, $A_1$ and $A_3$ differ only
in even-indexed coefficients (they share the odd choices, hence the same $O$), so they are
the endpoints of an even segment; both are Hurwitz by hypothesis, so by
Lemma~\ref{lem:seg} every matrix on this even segment is Hurwitz, whence its value
$\chi(\jj\omega)$ is nonzero for every $\omega\in\mathbb R$ (Lemma~\ref{lem:mik} (M1)).
But as the even coefficients vary along that even segment, $\chi(\jj\omega)$ traces, for
fixed $\omega\ge0$, the full horizontal segment
$\{x+\jj O^-(\omega):x\in[E^-(\omega),E^+(\omega)]\}$ (the value set of an even segment is
itself a one-dimensional interval, by the same affine-coordinate argument as in
Lemma~\ref{lem:rect}). Hence this entire bottom edge of $R(\omega)$ avoids $0$ for every
$\omega\ge0$. Identically: the top edge (from $A_4$ to $A_2$, an even segment with common
$O=O^+$), the left edge (from $A_1$ to $A_4$, an odd segment with common $E=E^-$), and the
right edge (from $A_3$ to $A_2$, an odd segment with common $E=E^+$) are origin-free for
all $\omega\ge0$. Thus
\begin{equation}\label{eq:edgefree}
0\notin\partial R(\omega)\qquad\text{for all }\omega\ge0,
\end{equation}
where $\partial R(\omega)$ is the boundary of the rectangle.

\emph{The origin is excluded for $\omega=0$ and for large $\omega$.} At $\omega=0$,
\eqref{eq:EO} gives $O_p(0)=0$ and $E_p(0)=a_0$ for every $p$; the corner $A_1$ being
Hurwitz forces its constant term to be positive (a Hurwitz monic real polynomial has
positive constant term, as shown in the proof of Lemma~\ref{lem:mik}), and $E^-(0)$ is the
minimal constant term over the box, attained by $A_1$, hence $E^-(0)>0$. So
\eqref{eq:inrect} fails at $\omega=0$ and $0\notin R(0)$. For large $\omega$, the dominant
term of $\chi_p(\jj\omega)$ is $(\jj\omega)^n$, of modulus $\omega^n\to\infty$ uniformly in
$p$ over the compact box (the remaining terms are bounded by
$C\omega^{n-1}$ with $C$ depending only on the fixed bounds $a_i^\pm$); hence there is
$\Omega>0$ with $|\chi_p(\jj\omega)|\ge \tfrac12\omega^n>0$ for all $\omega\ge\Omega$ and
all $p$, so $0\notin R(\omega)$ for $\omega\ge\Omega$.

\emph{Conclusion by a first-entry argument.} Suppose, for contradiction, that
$0\in R(\omega^*)$ for some $\omega^*>0$. The set
$S:=\{\omega\ge0:0\in R(\omega)\}$ is closed (the corner functions $E^\pm,O^\pm$ are
continuous in $\omega$, so the conditions \eqref{eq:inrect} are closed), is nonempty
($\omega^*\in S$), and is bounded above by $\Omega$; also $0\notin S$ since $0\notin R(0)$.
Let $\omega_0:=\inf S\in(0,\Omega]$. By closedness $\omega_0\in S$, i.e. $0\in
R(\omega_0)$. By minimality, $0\notin R(\omega)$ for $0\le\omega<\omega_0$. Now consider
the continuous family of rectangles $R(\omega)$. For $\omega<\omega_0$ the origin lies
outside $R(\omega)$; at $\omega=\omega_0$ it lies in $R(\omega_0)$. We claim
$0\in\partial R(\omega_0)$. Indeed, if $0$ were in the open interior of $R(\omega_0)$, then
since $E^-,E^+,O^-,O^+$ are continuous in $\omega$, the strict inequalities
$E^-(\omega_0)<0<E^+(\omega_0)$ and $O^-(\omega_0)<0<O^+(\omega_0)$ would persist for all
$\omega$ in a two-sided neighbourhood of $\omega_0$, giving $0\in R(\omega)$ for some
$\omega<\omega_0$ and contradicting minimality of $\omega_0$. Hence
$0\in\partial R(\omega_0)$, contradicting \eqref{eq:edgefree}. Therefore $S=\emptyset$,
i.e. $0\notin R(\omega)$ for all $\omega\ge0$.
\end{proof}

\section{Proof of Theorem~\ref{thm:main}}

\begin{proof}[Proof of Theorem~\ref{thm:main}]
\emph{Necessity.} If every matrix in $\mathcal A$ is Hurwitz, then in particular the four
matrices $A_1,A_2,A_3,A_4\in M_0\subset\mathcal A$ are Hurwitz.

\emph{Sufficiency.} Assume $A_1,A_2,A_3,A_4$ are Hurwitz. By Lemma~\ref{lem:zero},
$0\notin R(\omega)$ for all $\omega\ge0$; since $\chi_p(-\jj\omega)=\overline{\chi_p(\jj
\omega)}$, the same holds for $\omega<0$. Therefore, by Lemma~\ref{lem:det}, no member of
$\mathcal A$ has an eigenvalue on the imaginary axis:
\begin{equation}\label{eq:noimag}
\chi_p(\jj\omega)\ne0\qquad\text{for all }p\in\mathcal P,\ \omega\in\mathbb R .
\end{equation}

The coefficient box $\prod_i[a_i^-,a_i^+]$ is convex, hence path-connected, so the family
$\mathcal A$ is path-connected: for any $p\in\mathcal P$ there is a continuous path
$t\mapsto A(p_t)$ in $\mathcal A$, $t\in[0,1]$, with $A(p_0)=A_1$ (a Hurwitz member) and
$A(p_1)=A(p)$ (take $p_t$ the straight-line segment in coefficient space from the
coefficient vector of $A_1$ to that of $p$, which stays in the box by convexity). Along
this path the eigenvalues vary continuously (as in Step~2 of Lemma~\ref{lem:seg}). Define
\[
t_0:=\sup\{\tau\in[0,1]:A(p_t)\text{ is Hurwitz for all }t\in[0,\tau]\}.
\]
Since $A(p_0)=A_1$ is Hurwitz, $t_0$ is well defined and $t_0\ge0$. The matrix $A(p_{t_0})$
is a limit of Hurwitz matrices, so each of its eigenvalues has $\Real\lambda\le0$; by
\eqref{eq:noimag} none lies on the imaginary axis, so in fact $\Real\lambda<0$ for all of
them and $A(p_{t_0})$ is Hurwitz. If $t_0<1$, then by continuity of eigenvalues the matrix
$A(p_t)$ remains Hurwitz (all eigenvalues stay in the open left half-plane) for $t$ in a
neighbourhood of $t_0$, contradicting the definition of $t_0$. Hence $t_0=1$, so
$A(p)=A(p_1)$ is Hurwitz.

As $p\in\mathcal P$ was arbitrary, every matrix in $\mathcal A$ is Hurwitz.
\end{proof}

\begin{remark}\label{rem:gen}
The argument used Kharitonov's combinatorial coefficient theorem nowhere. Its skeleton is:
(i) a continuity/boundary-crossing principle stating that stability can be lost in a
connected matrix family only through an eigenvalue reaching the imaginary axis
(Theorem~\ref{thm:main}, sufficiency, and Step~2 of Lemma~\ref{lem:seg}); (ii) the
monotone-phase characterisation of Hurwitz matrices in terms of the determinant
$\det(\jj\omega I-A)$ (Lemma~\ref{lem:mik}); and (iii) a one-dimensional segment lemma
(Lemma~\ref{lem:seg}) showing that an edge of a matrix polytope with stable endpoints is
stable. The only place the special interval structure entered is Lemma~\ref{lem:rect}: for
companion matrices of an interval polynomial, the determinant value set on the imaginary
axis is an axis-parallel rectangle whose corners are the four distinguished vertices, so
that the four edges of the rectangle are exactly four edges of the matrix polytope and
Lemma~\ref{lem:seg} applies to them. For a general convex polytope of matrices the same
machinery (i)–(iii) yields the matrix Edge Theorem: the polytope is robustly stable iff
all of its one-dimensional edges are stable; the present theorem is the special case in
which checking the four Kharitonov edges suffices because the value set is a rectangle.
\end{remark}

\end{document}
